% ============================================================
%  Chapitre 7 — Réseaux Sémantiques
% ============================================================
\chapter{Réseaux Sémantiques}

\begin{tcolorbox}[notebox, title={Note}]
Ce chapitre est pris en charge par le binôme et sera complété avec les
implémentations finales. La structure ci-dessous définit le cadre théorique
et les éléments attendus.
\end{tcolorbox}

\section{Fondements théoriques}

\subsection{Définition}

Un \textbf{réseau sémantique} est un graphe orienté étiqueté dans lequel :
\begin{itemize}
  \item les \textbf{nœuds} représentent des concepts ou des individus ;
  \item les \textbf{arcs} représentent des relations entre ces entités.
\end{itemize}

\noindent Les relations les plus courantes sont :
\begin{description}
  \item[\code{is-a}] relation de spécialisation/généralisation
        (Rappeur \code{is-a} Artiste) ;
  \item[\code{instance-of}] relation d'appartenance d'un individu à une classe
        (soolking \code{instance-of} Rappeur) ;
  \item[propriétés] relations arbitraires entre nœuds
        (soolking \code{signe\_chez} universal).
\end{description}

\subsection{Héritage}

Le mécanisme d'\textbf{héritage} permet à un concept de récupérer les
propriétés de ses ancêtres dans la hiérarchie \code{is-a}.
L'héritage peut être :
\begin{itemize}
  \item \textbf{simple} : un seul parent ;
  \item \textbf{multiple} : plusieurs parents (peut engendrer des conflits) ;
  \item \textbf{avec exceptions} : une propriété héritée peut être redéfinie
        localement (ex : un pingouin \code{is-a} oiseau, mais ne vole pas).
\end{itemize}

\noindent L'héritage avec exceptions rend les réseaux sémantiques
\textbf{non-monotones} : ajouter un nœud ou un arc peut invalider une propriété
héritée.

\subsection{Exemple conceptuel — domaine musical}

\begin{figure}[H]
\centering
\begin{tikzpicture}[
  node distance=1.5cm and 2.5cm,
  concept/.append style={minimum width=3cm}
]
  % Concepts
  \node[concept] (Personne)  {Personne};
  \node[concept] (Artiste)   [below=of Personne]  {Artiste};
  \node[concept] (Rappeur)   [below left=of Artiste]  {Rappeur};
  \node[concept] (Chanteur)  [below right=of Artiste] {Chanteur};

  % Individus
  \node[concept, fill=yellow!20] (soolking) [below=of Rappeur] {soolking};
  \node[concept, fill=yellow!20] (universal) [right=3.5cm of Artiste] {universal};
  \node[concept, fill=green!10]  (Label) [above=of universal] {Label};
  \node[concept, fill=green!10]  (Major) [above right=of Label] {MajorLabel};

  % is-a arcs
  \draw[arrow, dashed] (Artiste) -- node[left, font=\footnotesize]{\code{is-a}} (Personne);
  \draw[arrow, dashed] (Rappeur) -- node[left, font=\footnotesize]{\code{is-a}} (Artiste);
  \draw[arrow, dashed] (Chanteur) -- node[right, font=\footnotesize]{\code{is-a}} (Artiste);
  \draw[arrow, dashed] (soolking) -- node[left, font=\footnotesize]{\code{inst.}} (Rappeur);
  \draw[arrow, dashed] (Major) -- node[right, font=\footnotesize]{\code{is-a}} (Label);
  \draw[arrow, dashed] (universal) -- node[right, font=\footnotesize]{\code{inst.}} (Major);

  % propriete
  \draw[arrow, color=red!70] (soolking) -- node[below, font=\footnotesize]
    {\code{signe\_chez}} (universal);
\end{tikzpicture}
\caption{Réseau sémantique — domaine musical (schéma)}
\label{fig:reseau-music}
\end{figure}

\section{Application — Industrie musicale}

\textit{[À compléter par le binôme : graphe des concepts musicaux,
requêtes d'héritage, propriétés héritées par les artistes,
démonstration des exceptions]}

\section{Application — Système judiciaire}

\textit{[À compléter par le binôme : graphe des acteurs judiciaires,
hiérarchie des infractions, héritage des peines,
conflits d'héritage pour les cas complexes]}

\section{Implémentation avec TweetyProject}

\textit{[À compléter par le binôme avec les extraits de code]}

\begin{lstlisting}[caption={Structure d'un réseau sémantique avec Tweety (exemple)}]
// A completer par le binome
// Exemple de structure attendue :
//
// SimpleGraph<SemanticNetworkNode> graph = new SimpleGraph<>();
// SemanticNetworkNode artiste = new SemanticNetworkNode("Artiste");
// SemanticNetworkNode rappeur = new SemanticNetworkNode("Rappeur");
// graph.add(artiste); graph.add(rappeur);
// graph.add(new DirectedEdge<>(rappeur, artiste, "is-a"));
\end{lstlisting}
